\documentclass[letterpaper,10pt,twocolumn]{article}
\usepackage[top=2.5cm,left=2.5cm,right=2.5cm]{geometry}%ajusta los margenes de las paginas
\usepackage[utf8]{inputenc}
\usepackage[spanish,english]{babel}
\usepackage{amsmath, amsthm, amsfonts, amssymb}%paquetes de la american mathematical society
\usepackage{latexsym}%paquete de simbolos matematicos
\usepackage{textcomp}%paquete con mas simbolos
\usepackage[x11names,table]{xcolor} %los nombres de colores disponibles estan /etc/X11/rgb.txt
\usepackage{colortbl}%para permitir el coloreado de tablas
\usepackage{float}%Tener mayor control sobre las figuras o flotantes
\usepackage{caption}%tiene un mayor control sobre las etiquetas de los flotantes
\captionsetup{labelfont=bf,figurewithin=section,tablewithin=section}
\usepackage{abstract}%paquete dedicado al control del abstract
\usepackage{appendix}
\usepackage[pdftex]{graphicx}
%\usepackage[dvips]{graphicx} %Para compilar con el comando latex, si las imagenes son .eps
\DeclareGraphicsExtensions{.png,.pdf,.jpg}
\graphicspath{{imagenes/}}%directorio donde se guardan las imagenes
\usepackage{hyperref}%para poder insertar direcciones web y crear hiperlinks dentro del
                     %documento
\spanishdecimal{.}%Utiliza el punto en vez de la coma en modo matemático
\numberwithin{equation}{section}
%\hyphenrules{nohyphenation}%Para que no parta las palabras de ninguna forma

\newcommand{\grado}{\textdegree}%no funciona en modo matematico
\newcommand{\espaciocaps}{\setlength{\abovecaptionskip}{2pt}
  \setlength{\belowcaptionskip}{0pt}}
\newcommand{\espacioitems}{\setlength{\itemsep}{0pt}}

%COMPILAR
% latex p1.tex 
% dvipdfm p1.dvi, para imagenes eps

%pdflatex p1.tex
%para pdftex 

%extraer imagenes de un pdf
%pdfimages archivo.pdf imaganes
%convert [options] input-file [options] output-file

\title{\bf Rayos X}
\author{ Facultad de Ciencias, Universidad Nacional Autónoma de México \\
  Laboratoio de Física Contemporánea \\
  Cosme Fulanito \\} 
\date{\today}

\begin{document}
\selectlanguage{spanish}
\renewcommand{\tablename}{Tabla}

\twocolumn[
\maketitle
\begin{onecolabstract}
  Se obtuvo el espectro característico del NaCl obteniendo 5 picos en los siguientes ángulos
  $12.4 \pm 0.5$, $14.0 \pm 0.5$, $27.6 \pm 1.0$, $31.2 \pm 0.6$ y $53.5 \pm 1.1$  que
  corresponde a las siguientes longitudes de ondas $1.2 \pm 1.5$, $1.4 \pm 1.3$, $1.3 \pm
  2.5$, $1.5\pm 1.4$ y $1.5 \pm 1.9$ respectivamente. 
\end{onecolabstract}
]

\section{Introducción}
\label{sec:introduccion}

Los rayos X fueron descubiertos en 1895 por el físico alemán Wilhelm Röetgen. Fueron llamados
X porque en ese momento no se sabía que tipo de radiación era, ahora se sabe que son fotones
altamente energéticos u ondas electromagnéticas con un longitud de onda del orden de 10 a
0.01 nm. En la actualidad se usan principalmente para penetrar materia sólida, afectar una
placa fotográfica, ionizar átomos  y provocar flourescencia de algunos compuestos químicos. 

El objetivo de esta práctica es obtener el espectro de rayos X característico del NaCl a 30
kV y 20 kV, y obtener las longitudes de onda correspondiente a cada pico mostrado en el
espectro.

\subsection{Principio físico}
\label{sec:principio}

Cuando un electrón es acelerado por cierto potencial y se dispara contra un átomo se puede
presentar uno de los siguientes dos casos: que electrón pase cerca pero sin impactar al átomo
o que el electrón se impacte con un electrón del átomo.

En el primer caso el electrón se desvía de su trayectoria debido a la fuerza de atracción
entre el núcleo (+) y el electrón (-), y se frena perdiendo energía que es emitida como
radiación, este tipo de radiación es llamada Bremsstrahlung\footnote{Del idioma alemán, que
  significa radiación de frenado.}. En el espectro de rayos X es la parte continua y suave,
ver fig. \ref{fig:espectro}.Si el \cite{gautreau}. Dependiendo del potencial usado este
espectro continuo tiene una longitud de onda mínima desde el cuál comienza, la cual es igual
a,

\begin{equation}
  \label{eq:lmin}
  \lambda_{min} = \frac{hc}{eV}
\end{equation}

que es cuando toda la energía cinética del electrón es convertida en radiación, esto es que
el electrón es frenado en seco.

En el segundo caso la energía cinética del electrón es tal que logra desprender al electrón
atómico dejando un hueco que inmediatamente es ocupado por un electrón del mismo átomo en un
nivel de mayor energía que al caer emite radiación\footnote{O un fotón.}, de un valor igual a
la diferencia de energía entre los dos niveles involucrados. En el espectro de rayos X
corresponde a los picos de la gráfica, ver fig. \ref{fig:espectro}. Existe una notación del
tipo $A_{i}$, para designar entre que niveles\footnote{En espectroscopía también se les suele
  llamar capas.} ocurre la emisión del fotón, donde $A$ indica la capa en la que fue arrancado
el electrón y el subíndice $i$ indica el nivel de donde cayó el electrón que lleno el hueco,
es decir si hueco es dejado en el nivel 1 se le asocia la letra $K$\footnote{En la notación
  espectroscópica el número cuántico $n=1$ se designa por $K$, $n=2$ por $L$, $n=3$ por $M$
  y así sucesivamente.} y si es llenado por un electrón del siguiente nivel superior, $n=2$,
se le asigna la letra griega $\alpha$ en cambio si se llena con un electrón que estaba dos
niveles arriba se le asigna la letra $\beta$ y así sucesivamente, por lo que para un fotón
emitido al caer del nivel 5 al nivel 2 se le asignaría como  $L_{\gamma}$

\begin{figure}[H]
  \espaciocaps
  \centering
  \includegraphics[width=7cm,height=6cm]{espectro}
  \caption{Espectro típico de rayos X generados por un ánodo de Mo a una diferencia de
    potencial de 35 KeV.\cite{serway}}
  \label{fig:espectro}
\end{figure}

\subsection{Producción de rayos X}
\label{sec:produccion}

Una forma común de producir rayos X es por medio de un tubo de rayos catódicos como el de la
fig. \ref{fig:tubo}. El cátodo, por medio del filamento, es calentado por lo que los átomos
del cátodo comienzan a vibrar provocando que la energía para desprender electrones
superficiales disminuya, a esta forma de producir electrones se le llama emisión
termoiónica\cite{beisser} PAG. 67(81). La gran diferencia de potencial aplicada entre el
ánodo y el cátodo genera un haz de electrones que va a chocar contra el ánodo ocurriendo el
proceso descrito en \ref{sec:principio}, al ser bombardeado constantemente puede llegar
a calentarse demasiado por eso se utilizan metales con un punto de fusión elevado entre los 
más comunes son tungsteno, molibdeno o cobre.\cite{beisser}

\begin{figure}[H]
  \centering
  \includegraphics[width=7cm]{tubo}
  \caption{Tubo de rayos catódicos.\cite{serway}}
  \label{fig:tubo}
\end{figure}

\subsection{Detección de rayos X}
\label{sec:deteccion}

Para detectar la radiación electromagnética\footnote{Los rayos X interactuan con la materia
  principalemente mediante tres procesos: absorción fotoeléctrica, dispersión de Compton y
  producción de pares.\cite{krane}} existen diversos de aparatos como son: cámaras 
de burbujas, centelleadores, cámaras de ionización, contadores proporcionales, detectores
Geiger-Müller, entre otros. Estos últimos tres funcionan por medio de una diferencia de
potencial aplicada a un gas, la diferencia entre los tres depende del voltaje aplicado entre
el ánodo y cátodo porqué de esto depende los procesos que ocurren dentro de cada uno de
ellos, fig. \ref{fig:regiones}.

\begin{figure}[H]
  \espaciocaps
  \centering
  \includegraphics[width=8cm,height=6cm]{regiones}
  \caption{Las diferentes regiones de trabajo de ionización que dependen del voltaje aplicado
    y la corriente producida por la creación de iones dentro del gas.} 
  \label{fig:regiones}
\end{figure}

El funcionamiento del tubo GM es el siguiente: la radiación incidente pasa a través de una
ventana delgada, fig. \ref{fig:geiger}, y entra al tubo donde ioniza el gas dejando pares de
iones y electrones a su paso estos, debido a la gran diferencial de potencial entre el cátodo
(el interior del tubo) y el ánodo (alambre) los iones positivos son atraídos al cátodo
mientras que los electrones son atraídos al ánodo. En el camino al ánodo los electrones van
acelerándose y chocando con otros átomos del gas a los que también ioniza generando más pares
de iones, cada electrón puede llegar a ionizar hasta dos átomos una corriente cientos de
veces mayor que en un contador proporcional. Al ser colectados por el ánodo los electrones se
genera una señal negativa por otro lado los iones positivos al ser más masivos tardan mucho
más en llegar al cátodo y cuando lo hacen generan un señal positiva, el tiempo transcurrido
entre estas dos señales es llamado tiempo muerto y durante él cualquier otra ionización
ocurrida en el tubo GM no es registrada\footnote{Por lo que un tubo GM no puede determinar
  eventos más cercanos que $10^{4}$ por segundo.}. Ocurre que también al impactarse los iones 
positivos contra los átomos del cátodo se vuelven a desprender electrones que son atraídos al
ánodo sin embargo esta es una situación que no se desea ya que genera corrientes espurias,
para evitar esto aparte del gas de ionización se utiliza un gas llamado de enfriamiento para
que estos nuevos electrones en lugar de generar una nueva corriente pierdan sus energía en
deshacer las moléculas del gas de enfriamiento\footnote{Puede ser orgánico ethanol o
  inorgánico como el bromo o el cloro.}. También tomando en cuenta el tiempo de este proceso
de 'enfriamiento' el tiempo entre cada ionización debe ser de alrededor de $40 \mu
s$.\cite{littlefield} 

\begin{figure}[H]
  \espaciocaps
  \centering
  \includegraphics[width=8cm]{geiger}
  \caption{Esquema típico del tubo Geiger-Müller.\cite{serway}}
  \label{fig:geiger}
\end{figure}

\subsection{Ley de Bragg}
\label{sec:bragg}

Una onda electromagnética en la región del visible al pasar a través de una rejilla de
difracción es difractada obteniendo mínimos y máximos en el lugar donde la onda tuvo
interferencia destructiva o constructiva, respectivamente. William Bragg, padre e hijo, en
Inglaterra y Max von Laue en Alemania sugirieron que un cristal\footnote{Un arreglo periódico
  de átomos.} pudiese actuar como una rejilla de difracción tridimensional para radiación de
cierta longitud de onda, de hecho para una $\lambda \le 2d$, donde $d$ es el espacio
entre planos sucesivos del cristal, ver fig. \ref{fig:bragg}.\cite{serway}

En la fig. \ref{fig:bragg} el rayo 1\footnote{Los rayos 1 y 2 evidentemente son rayos X.}, al
chocar con un átomo del plano A del cristal es reflejado especularmente a cierto ángulo
$\theta$ mientras que el rayo 2 es reflejado, al mismo ángulo, sobre un átomo en un plano
sucesivo más abajo (el plano B), si la diferencia de camino $ABC$, entre los dos rayos, es
igual a un número entero de longitudes de onda a la salida se obtiene interferencia
constructiva la cual puede ser detectada. Por lo que se tiene la siguiente
relación\cite{serway}: 

\begin{equation}
  \label{eq:bragg}
  n\lambda = AB + BC = 2d \sen\theta. 
\end{equation}

\begin{figure}[H]
  \espaciocaps
  \centering
  \includegraphics[width=8cm]{bragg}
  \caption{Ley de Bragg.\cite{serway}}
  \label{fig:bragg}
\end{figure}

\section{Desarrollo experimental}
\label{sec:desarr}

En la figura \ref{fig:dispositivo} se encuentra el diagrama esquemático del experimento. Para
la primera etapa, la de alto voltaje, se fijo un voltaje de salida de 580 V. En el diagrama el
llamado contador\footnote{Ratemeter.} se colocó a una escala máxima de 500 cuentas por 
segundo. El dispositivo de control es el encargado tanto de hacer mover el motor, recibir la
señal que llega del contador y enviarla a la computadora como también enviar a la computadora
la posición del tubo GM, que el doble del ángulo, al cual el rayo X es difractado. 

\begin{figure}[H]
  \centering
  \includegraphics[width=8cm]{dispositivo}
  \caption{Diagrama esquemático.}
  \label{fig:dispositivo}
\end{figure}

El espectrómetro\footnote{Un espectrómetro típico se puede ver en la
  fig. \ref{espectrometro}.} de rayos X utilizado fue el Tel-X-Ometer Tel.580\footnote{De la
  marca Teletron.}, además se ocuparon los siguientes accesorios: el colimador primario de
rendija vertical (582.001), el colimador de rendija vertical de 3 mm (562.016), el colimador
de rendija vertical de 1 mm (562.015) y el detector Geiger-Müller. En el brazo movible del
Tel.580 se colocó el tubo GM en la ranura 26, enfrente del tubo en la ranura 18 se colocó el
colimador de 1 mm y el colimador de 3 mm en la ranura 13 mientras que el  colimador primario
de rendija vertical fue colocado a la salida del tubo de rayos catódicos. El cristal de NaCl
fue colocado en el portacristales del espectrómetro cuidando que la cara menos brillante
apuntara hacia el tubo de rayos X. 

\begin{figure}[H]
  \centering
  \includegraphics[width=8cm]{espectrometro}
  \caption{Espectrómetro de rayos X.\cite{serway}}
  \label{espectrometro}
\end{figure}

El rango de trabajo fue de 12\grado a 120\grado en la escala de $2\theta$, en la computadora
se utilizo un retardo de 75 ms y se hicieron 7 mediciones para 30 kV y 5 para 20 kV. 

\section{Resultados}
\label{sec:resultados}

Los resultados de las mediciones mostraron claramente tres picos principales. También
aparecieron picos secundarios muy cercanos a los primeros dos picos principales, ver
fig. \ref{fig:30kv} y \ref{fig:20kv}. Los resultados completos se pueden ver en la tabla
\ref{tab:resultados}, para la distancia interplanar se utilizo la dada en el manual del
espectrómetro \cite{teletron}, que es igual a $d = 0.282$ nm.

\begin{table}[H] 
\centering
\rowcolors{3}{gray!30}{}
\begin{tabular}{|c|c|c|c|c|c|} 
\hline 
\rowcolor{LightBlue2} $2\theta$ & $\theta$ & $\sin\theta$ & $2d$ & $n\lambda$ & $n$ \\  
\rowcolor{LightBlue2} (\grado) & (\grado) &  & (nm) & ($\AA$) &  \\  
\hline 
24.8 & 12.4 & 0.22 & 0.564 & 1.2 & 1 \\ 
 $\pm$ 0.5 & & & & $\pm$ 1.5 & \\ 
\hline 
27.9 & 14.0 & 0.24 & 0.564 & 1.4 & 1 \\ 
 $\pm$ 0.5 & & & &  $\pm$ 1.3 & \\ 
\hline 
55.2 & 27.6 & 0.46 & 0.564 & 2.6 &  2 \\ 
 $\pm$ 1.0 & & & & $\pm$ 2.5 & \\ 
\hline 
62.5 & 31.2 & 0.52 & 0.564 & 2.9 & 2 \\ 
 $\pm$ 0.6 & & & & $\pm$ 1.4 & \\ 
\hline 
106.9 & 53.5 & 0.80 & 0.564 & 4.5 & 3 \\ 
 $\pm$ 1.1 & & & & $\pm$ 1.9 & \\ 
\hline 
\end{tabular} 
\caption{Resultados del experimento.} 
\label{tab:resultados}
\end{table} 

Las gráficas de las 7 mediciones de datos para 30 kV están en la fig. \ref{fig:30kv}
mientras que los de las 5 mediciones para 20 kV están en la fig. \ref{fig:20kv}.

\begin{figure}[H]
  \centering
  \includegraphics[width=8cm,trim= 0.7cm 0.3cm 1.7cm 0.8cm,clip]{30kv}
  \caption{Espectro del NaCl a 30 kV.}
  \label{fig:30kv}
\end{figure}

\begin{figure}[H]
  \centering
  \includegraphics[width=8cm,trim= 0.7cm 0.3cm 1.7cm 0.8cm,clip]{20kv}
  \caption{Espectro del NaCl a 20 kV.}
  \label{fig:20kv}
\end{figure}

\section{Conclusiones}
\label{sec:conclusiones}

Con las mediciones realizadas se puede comprobar que efectivamente los rayos X generados por
el espectrómetro son del orden la distancia interatómica del NaCl, sin embargo dadas las
incertidumbres no se puede decir más cosas con mayor certeza. El hecho de que no hayan sido
tan buenas las mediciones en parte se debe a que el tiempo de retardo, que ahora creemos, debe
ser mayor a 75 ms\footnote{Aunque esto generaría que cada medición se llevará arriba de 10
  minutos} y al sistema de automatización que algunas veces se llegaba a atorar, lo que
provocaba volver a realizar la medición.

Para obtener las $\lambda$ sólo fue necesario obtener los ángulos en los que se presentaban
los picos, su altura sólo nos indica cuanta radiación paso a través del GM.

\bibliographystyle{unsrt}
\bibliography{referencias.bib}
\end{document}
